\section{Introduction}
\label{sec:introduction}

AI-assisted programming tools have transitioned from experimental to standard practice across the software industry. As of 2025, 85\% of developers regularly use AI tools for coding and software design, and 51\% of professional developers report using them daily~\cite{modall2025}. Approximately 46\% of newly written code is now AI-assisted, a figure projected to reach 60\% by the end of 2026~\cite{modall2025}. GitHub Copilot alone has accumulated 20 million cumulative users and is deployed by 90\% of Fortune 100 companies, while the AI coding tools market reached \$7.37 billion in 2025~\cite{modall2025}. Systems such as Copilot, Devin, and Cursor automate tasks ranging from code completion to full pull request authorship, with minimal human oversight at each step.

This adoption is accompanied by documented quality concerns. Developer trust in AI output has declined sharply: only 29\% of developers reported confidence in AI-generated code in 2025, down from above 70\% in 2023, and 45\% say that debugging AI-generated code takes more time than writing equivalent code manually~\cite{modall2025}. Industry analyses report 2.74 times more vulnerabilities in AI-generated code relative to human-authored code, with 45\% of AI code samples failing security tests outright~\cite{modall2025}. Prior research confirms that AI-generated code exhibits higher defect rates, yet leaves the broader picture incomplete in four ways. First, most studies are restricted to Python or Python and Java, leaving cross-language patterns and the moderating role of type discipline untested at scale. Second, each study typically examines a single quality dimension -- stylistic violations, defect rates, or security findings -- rather than applying a unified framework. Third, statistical approaches vary from descriptive-only reporting to parametric tests on non-normal distributions, with effect sizes absent throughout. Fourth, no existing work evaluates multiple production-level LLMs jointly on the same benchmark, making it unclear whether observed patterns reflect a particular model or AI-generated code more broadly.

We address this through a large-scale empirical analysis of the CodeMirage benchmark~\cite{codemirage2024}, a multilingual dataset comprising approximately 110,000 files spanning 10 programming languages and generated by 10 production-level LLMs alongside a human-authored baseline. The LLM pool includes both closed-weight models (Claude-3.5-Haiku, GPT-4o-mini, GPT-o3-mini, Gemini-2.0-Flash, Gemini-2.0-Flash-Thinking, Gemini-2.0-Pro) and open-weight models (DeepSeek-V3, DeepSeek-R1, Llama-3.3-70B, Qwen-2.5-Coder-32B), enabling a direct comparison of model families alongside the AI-versus-human analysis. We apply language-specific linters, structural complexity measurement, and security scanning uniformly across all languages, with all comparisons grounded in nonparametric statistical tests and effect sizes, stratified by type discipline to examine whether static type systems moderate AI-associated defect risk. This study examines three research questions:

\vspace{0.4em}
\noindent \textbf{RQ1:} \textit{What structural and stylistic quality differences exist between AI-generated and human-written code?}

\vspace{0.2em}
\noindent \textbf{RQ2:} \textit{How do security vulnerability profiles differ between AI-generated and human-written code?}

\vspace{0.2em}
\noindent \textbf{RQ3:} \textit{Do open-weight LLMs produce code with measurably different quality and security profiles compared to closed-weight LLMs when evaluated on the same benchmark dataset?}
\vspace{0.4em}

\noindent The main contributions of this work are as follows:

\begin{enumerate}

\item \textbf{Large-scale multi-language empirical comparison across 10 languages and 10 LLMs.} We analyse the CodeMirage benchmark using a unified framework covering stylistic, structural, and security quality dimensions, providing systematic evidence of how AI authorship and model family affect code quality across language families.

\item \textbf{Evidence that type discipline moderates AI-associated quality gaps.} Statically-typed languages (C, C\#, Java, Go) exhibit a mean 1.66-fold violation density gap between AI-generated and human-authored code, compared to a 2.34-fold gap in dynamically-typed languages (Python, JavaScript, PHP, Ruby). A chi-square test confirms this pattern ($p < 0.001$), suggesting that compiler-enforced type checking intercepts a class of errors that otherwise persist in dynamic languages.

\item \textbf{Characterisation of AI-specific security vulnerability patterns.} AI-generated code carries 2.75 times higher vulnerability density than human-authored code (3.71 vs.\ 1.35 findings per KLOC), with 69.0\% of AI-generated files containing at least one security finding compared to 44.0\% of human-authored files. CWE category distributions differ significantly by code origin ($\chi^2(4) = 72.0$, $p < 0.001$, Cram\'{e}r's $V = 0.30$), with AI-generated code concentrated in injection and secrets-related categories.

\item \textbf{Finding that structural complexity metrics do not distinguish AI from human code.} Across 19 stylometry and complexity metrics, only two reach statistical significance after Bonferroni correction, both with negligible effect sizes ($|\delta| \leq 0.14$). Code review frameworks that rely on structural complexity thresholds alone are therefore unlikely to surface AI-specific defects.

\end{enumerate}
